% Created 2024-01-31 Wed 16:35
% Intended LaTeX compiler: pdflatex
\documentclass[11pt]{article}
\usepackage[utf8]{inputenc}
\usepackage[T1]{fontenc}
\usepackage{graphicx}
\usepackage{longtable}
\usepackage{wrapfig}
\usepackage{rotating}
\usepackage[normalem]{ulem}
\usepackage{amsmath}
\usepackage{amssymb}
\usepackage{capt-of}
\usepackage{hyperref}
\date{\today}
\title{}
\hypersetup{
 pdfauthor={},
 pdftitle={},
 pdfkeywords={},
 pdfsubject={},
 pdfcreator={Emacs 29.1 (Org mode 9.6.17)}, 
 pdflang={English}}
\usepackage{changes}
\begin{document}

\tableofcontents

\section{Zero-covariation cases}
\label{sec:orgde51e43}

\textbf{pp. 2--3}: Cheng discusses several examples of causal judgments when
covariation between a cause and an effect is zero. What should one
conclude if \(P(e|i) - P(e|\bar i)=0\)? There is a problem, for example,
if another cause \(j\) is always present, because then one observes
\(P(e|ji) - P(e|j\bar i)=0\) and cannot conclude that \(i\) is not causing
\(e\) in the absence of tests with \(j\) alone.

Empirically, people judge \(P(e|i)=P(e|\bar i)\simeq 1\) as \(i\) having
uncertain causal status, while \(P(e|i)=P(e|\bar i)<1\) leads to
judgments that \(i\) is not a cause of \(e\).

Cheng also notes that inhibitory causes cannot be detected if the
effect does not occur with some frequency in the absence of the cause
(well known in associative learning for conditioned inhibition). In
this case one has \(P(e|i)=P(e|\bar i)\simeq 0\).

\textbf{p4} Aims to contrast PowerPC model and RW model / delta-rule.

\textbf{p6} Mentions learned irrelevance but explanation is not super clear,
 is it mentioned later?

\section{Main-Effect Contrast and Generative Causal Power (p6 ff.)}
\label{sec:org138f5ae}

\begin{itemize}
\item \(p_x\) is the Pr that \(x\) causes \(e\) when \(x\) is present
(causal power of \(x\)).
\item \(p_x\) is not the same as \(P(e|x)\) (Pr that \(e\) occurs when \(x\) is
present).
\item \(P(e|x)\) can be estimated from observations, while \(p_x\) must be
estimated using observations and other assumptions.
\item \(p_x = P(e|x)\) only when no other causes are present or exist.
\item Let \(i\) be a candidate cause for which we want to estimate \(p_i\),
and \(a\) all alternative causes.
\end{itemize}

In summary:
\begin{itemize}
\item when \(i\) occurs, it produces \(e\) with Pr \(p_i\).
\item when \(a\) occurs, it produces \(e\) with Pr \(p_a\).
\item nothing else influences \(e\).
\item \(i\) and \(a\) influence \(e\) independently.
\item causal powers \(p_i\) and \(p_a\) are independent of how often \(i\) and
\(a\) occur. That is:
\begin{displaymath}
  P(e) = P(i)p_i + P(a)p_a - P(i)p_i\times P(a)p_a
\end{displaymath}
\item with \(P(i)=1\) and \(P(i)=0\) one gets \(P(e|i)-P(e|\bar
  i)=(1-P(a|i)p_a)p_i+(P(a|i)-P(a|\bar i))p_a\)
\item When \(i\) and \(a\) are independent, \(P(a|i)=P(a|\bar i)\) and we get
\begin{equation}
  \label{eq:pi}
  p_i = \frac{ P(e|i) - P(e|\bar i) }{ 1 - P(a)p_a }
\end{equation}
\item Cheng points out that, although \(p_a\) is unobservable, the product
\(P(a)p_a\) can be estimated from observations of when \(a\) occurs
without \(i\). \textbf{This product looks the same as \(V_a\) in associative
learning.}
\item The formal statement of the previous point is \(P(a)p_a=P(e|\bar i)\),
leading to
\begin{equation}
  \label{eq:pi}
  p_i = \frac{ P(e|i) - P(e|\bar i) }{ 1 - P(e|\bar i) }
\end{equation}
\end{itemize}

\section{Main-Effect Contrast and Preventive Causal Power}
\label{sec:orgdda9ae0}

\begin{itemize}
\item Now \(p_i\) is the probability that \(i\) prevents \(e\).
\item In this case \(e\) occurs if \(a\) causes it and \(i\) does \emph{not} prevent
it:
\begin{equation}
  \label{eq:pipreventive}
  P(e) = P(a) p_a ( 1 - P(i)p_i )
\end{equation}
\item When \(a\) and \(i\) are independent:
\begin{equation}
  p_i = \frac{ - (P(e|i) - P(e|\bar i)) }{ P(e|\bar i) }
\end{equation}
\end{itemize}
\end{document}
